\chapter{Diseño eléctrico y seguridad del sistema aislado}
\section{Arquitectura, transformación y límites}
La generación opera inicialmente aislada del SEIN. Tres generadores de 250 kVA se conectan a 380 V, con interruptor por grupo y barra principal; el respaldo diésel tiene tres módulos de 200 kW. Se limita la potencia total simultánea de generación a 600 kW. Un transformador de 800 kVA, 10/0.38 kV, eleva hacia una red primaria académica de 8 km y seis subestaciones de distribución: tres de 160 kVA, dos de 100 y una de 50 kVA, total 730 kVA. Las potencias de placa de distribución no se suman como demanda coincidente.

El grupo de conexión propuesto mantiene referencia AT/BT: devanado de 10 kV en delta y devanado de 380 V en estrella con neutro accesible, con desplazamiento horario por especificar con el fabricante. Invertir el flujo de potencia no permite renombrar arbitrariamente el grupo vectorial. El transformador auxiliar de 20 kVA es 380/380--220 V, con estrella secundaria y neutro propio para utilización.
\begin{figure}[htbp]\centering\includegraphics[width=.98\linewidth]{figuras/unifilar-r03.png}\caption{Unifilar académico del sistema aislado. Sin fuente SEIN; impedancias, interruptores y neutros requieren coordinación definitiva.}\end{figure}
\input{generado/balance}
Las pérdidas del cuadro se descuentan una sola vez, usando factor de potencia de transporte 0.90 como hipótesis conservadora; la demanda doméstica conserva 1.0 de N21. Resultan \dato{Pusers} kW en usuarios y margen de \dato{Margin} kW sobre el pico. Arranques, temperatura, desequilibrio o caída adicional pueden consumir ese margen. El balance no demuestra cumplimiento de calidad.

\section{Conductores y caída de tensión}
La corriente nominal por generador es \dato{Iunit} A y la suma de placas, \dato{Ibus} A. El transformador de 800 kVA tiene corriente nominal BT de \dato{Itrafo} A. Se estudia interruptor de 630 A por grupo y de 1600 A en barra, con capacidad de interrupción candidata de 25 kA. Para equipos marcados al 80\% de carga continua, 630 y 1600 A permiten 504 y 1280 A respectivamente. La selección final debe considerar temperatura, fabricante, modo de instalación y coordinación.

La regla CNE Utilización 050-104 distingue corriente continua, marcado del equipo y tablas de ampacidad. Las condiciones de las tablas de enterramiento D8A/D8B son específicas: no corresponde multiplicar indiscriminadamente toda ampacidad por 0.8. Como referencia, para cobre de 240 \(\mathrm{mm^2}\) el detalle 4 consultado indica 446 A en D8A y 442 A en D8B, y otros detalles dan valores diferentes. Se proponen dos conductores de cobre de 240 \(\mathrm{mm^2}\) por fase y generador, longitud 25 m. Cada uno conduce aproximadamente 190 A, pero su ampacidad admisible debe verificarse con disposición, agrupamiento, temperatura, resistividad térmica y reparto entre paralelos. Se requiere igualdad de material, sección y longitud, conexiones equivalentes y PE dimensionado según el sistema.

La resistencia de estudio por conductor es 0.095 ohm/km; la caída del ramal de generación es aproximadamente \dato{DropCable}\%. Los 450 m presupuestados corresponden a \(3\) grupos por \(3\) fases por \(2\) paralelos por \(25\) m; no incluyen ficticiamente un neutro de potencia en cada fase. Las conexiones de respaldo se consideran con cuatro paralelos de cobre de 400 \(\mathrm{mm^2}\) por fase a lo largo de 25 m, sujeto a revisión térmica.

Para la línea primaria se propone AAAC de 95 \(\mathrm{mm^2}\), resistencia a 20 grados de 0.35 ohm/km y coeficiente térmico 0.0036/\(\!^\circ\mathrm{C}\). A 40 grados resulta aproximadamente 0.3752 ohm/km. La reactancia 0.38 ohm/km es una hipótesis de disposición; deben fijarse geometría y flechas. A plena carga la caída es \dato{DropMT}\%. El recorrido de 8 km es un corredor aproximado para cubrir las localidades, no una línea levantada ni un trazado aprobado. El diseño mecánico de postes, vanos y distancias CNE debe desarrollarse con topografía real.

Se evalúa una rama BT de 25 kW, 150 m, 380 V y conductor autoportante de aluminio 3 por 35 más 16/25 \(\mathrm{mm^2}\), según cuadro de N21: \(R_{40}=0.930\) ohm/km, \(X=0.103\) ohm/km y corriente de referencia 129 A. Resulta corriente aproximada de 38 A y caída \dato{DropBT}\%. La norma fija sección mínima de fases de 25 \(\mathrm{mm^2}\); no se arrastra la antigua selección de 16 como conductor general de fases. Su numeral 4.5 contempla 7.5\% máximo BT para SER bajo su ámbito; ello no reemplaza los criterios de utilización ni demuestra tensión correcta en todos los usuarios. La regulación del transformador, reparto por fase, longitud real, arranque y calidad deben analizarse conjuntamente. La pérdida total BT de 3\% es una hipótesis del balance, no el resultado de todos los circuitos individualizados.

\section{Cortocircuito y protección de una red débil}
Se utiliza reactancia subtransitoria supuesta de 0.20 p.u. para cada generador hidráulico, 0.15 p.u. para respaldo y \(u_k=6\%\) para T1. Se calcula \(I''\simeq I_n/X''\) y se combina la impedancia equivalente con el transformador. El resultado inicial simétrico es \dato{IKH} kA para hidro, \dato{IKB} kA para respaldo y una cota simplificada de \dato{IKmax} kA para fuentes simultáneas. La contribución hidroeléctrica referida al lado 10 kV es \dato{IKMT} A. Estos datos no son valores de placa; se requiere impedancia de secuencia, excitación, motores, fuente diésel, posición y tipo de falla.

IEC 60909-0:2026, edición 3 (N25), sustituyó la edición de 2016 según la ficha oficial. No se obtuvo el texto completo, por lo que no se declara que este cálculo simplificado cumpla íntegramente el método IEC. El Icu candidato de 25 kA debe revisarse con corriente de cierre, componente continua, relación X/R, capacidad térmica y dinámicas. Los cortocircuitos monofásicos y extremos de red no pueden inferirse del máximo trifásico de barra.

La corriente de falla de un generador disminuye con el tiempo: con \(X_d=1.8\) p.u. supuesto, el agregado hidráulico puede aproximarse a 633 A sostenidos, inferior a la corriente de carga plena. \textbf{Una protección de sobrecorriente 50/51 por sí sola puede no detectar una falla remota}. Se prevén protección diferencial de generador y transformador (87G/87T), sobrecorriente restringida o controlada por tensión (51V), fallas a tierra de acuerdo con neutros (64G/59N/67N), potencia inversa (32), pérdida de excitación (40), secuencia negativa (46), tensión y frecuencia (27/59/81), y comprobación de sincronismo (25) entre grupos internos. Los ajustes necesitan curvas de decremento, capacidades térmicas, transformadores de medida y estudio de selectividad; no se publican ajustes inventados.
\input{generado/fallas_remotas}
El cuadro agrega impedancia de 8 km de línea, transformación SED con \(u_k=4\%\) supuesto y rama BT de 150 m. Se compara operación de tres grupos, un solo grupo y un límite sostenido de un grupo, sin asumir que este último sea una falla estacionaria garantizada con excitación real. La caída de corriente confirma la necesidad de sensibilidad y selectividad. No se derivan de este cuadro corrientes fase--tierra, impedancia homopolar ni ajustes definitivos.

\section{Neutros y puesta a tierra}
Se propone estudiar puesta a tierra del neutro de generación mediante impedancia coordinada y un único esquema activo compatible con grupos en paralelo. La delta de 10 kV no aporta un neutro: si se adopta puesta a tierra de esa red por impedancia será necesario transformador zigzag y resistencia dimensionados. Las subestaciones de distribución y el auxiliar tendrán sus neutros BT y conductores de protección definidos según el sistema de utilización. No se conectan resistencias de neutro como si fueran conductores PEN de viviendas.

La malla conceptual de casa de máquinas es 20 por 12 m, paso 4 m, cobre de 50 \(\mathrm{mm^2}\), longitud horizontal 152 m y ocho electrodos de 2.4 m. La sección no se certifica sin corriente a tierra y duración de despeje. El estimador de gabinete \(R_g\simeq\rho/(4\sqrt{A/\pi})+\rho/L\) sirve para sensibilidad, sin acreditar el aporte de varillas ni efectos de capas del suelo.
\input{generado/tierra}
El GPR es elevación de potencial de la malla y \textbf{no equivale a tensión de contacto}. CNE Suministro 036 exige considerar paso, contacto y transferencia; su referencia de 25 ohm para un electrodo individual no es un criterio global suficiente de seguridad. Tampoco existe aquí una resistencia universal de 5 ohm que cierre la verificación. Se requiere medir resistividad por estratos, definir reparto de corriente, tiempo de despeje, superficie y cercos, y calcular tensiones accesibles con metodología reconocida. La ficha IEEE 80-2013 consultada indica condición reservada/inactiva; se debe seleccionar y adquirir una edición contractual reconocida, sin presentarla como texto completo vigente consultado.
\begin{figure}[htbp]\centering\includegraphics[width=.95\linewidth]{figuras/puesta-tierra-r03.png}\caption{Malla y separación conceptual de funciones de neutro. Dimensiones de estudio; verificación de paso/contacto pendiente.}\end{figure}

\section{Espacios y seguridad}
CNE Suministro 125 y tablas correspondientes definen espacios según tensión y condición de exposición. Para BT 151--600 V, los casos de la tabla 125-1 consultada exigen 900, 1070 o 1200 mm; no se aplica una sola distancia a cualquier tablero. La planta académica reserva 1.5 m frontal, y la media tensión debe contrastarse con su tabla aplicable y equipo real. Se requieren segregación, acceso restringido, barreras, señalización, bloqueo y consignación, extinción compatible, control de incendio de combustible, salidas y plan de rescate. Distancias aéreas de suministro y zonas de trabajo se verifican en el proyecto de distribución, considerando relieve y cruces reales.

\section{Adaptación eventual al SEIN entre 2046 y 2076}
La memoria oficial de Electro Puno 2023 incluye transformación 10/22.9 kV y redes regionales de 10, 22.9 y 60 kV (H14). No acredita un punto disponible ni una conexión aprobada en Limbani. Se adopta 10 kV en el diseño actual por el antecedente local; no se elige automáticamente 22.9 kV para simular una futura interconexión.

Como previsión se reserva espacio para una celda adicional, canalizaciones de comunicaciones, medición y acceso a transformación exterior. Puede estudiarse equipamiento con aislamiento de clase 24 kV cuando su costo y funcionamiento a 10 kV resulten adecuados, pero elevar la tensión futura exigiría verificar o sustituir transformadores, aislamiento, conductores, distancias y protecciones. Las decisiones se tomarán con la red y normativa existentes en aquel momento.

Una interconexión requeriría estudio de conexión y autorización de la concesionaria, flujos de carga, cortocircuito, estabilidad, sincronismo, puesta a tierra, calidad, protección de interfaz y antiisla, medición comercial, comunicaciones, actualización ambiental, servidumbres y concesiones pertinentes. El horizonte 20--50 años es orientativo; no altera la operación aislada ni se utiliza como garantía de crecimiento o rentabilidad.
